\begin{table}[ht]
\centering
\caption{DSR Research Elements Mapping}
\label{tab:dsr-elements}
\begin{tabularx}{\textwidth}{|l|X|}
\hline
\textbf{DSR Element} & \textbf{Application in this Thesis} \\ \hline
\textbf{Context} & Cloud-native environments using Kubernetes; DevOps and SRE practitioners aiming for reliability under resource constraints. \\ \hline
\textbf{Problem} & Lack of unified reference architectures for multimodal observability leads to elevated Mean Time to Diagnose ($MTTDiag$), diagnostic inaccuracy under cascading faults, and unquantified resource overhead (Wicked Problem). \\ \hline
\textbf{Theoretical Conjectures} & Grounding observability pipelines in Control Principles (\texttt{CP1}, \texttt{CP2}) and Contingency Heuristics (\texttt{ConH1}--\texttt{ConH3}) significantly reduces $MTTDiag$ ($H_{1,1}$), increases root-cause accuracy ($H_{1,2}$), and confines resource overhead within a bounded $15\%$ envelope ($H_{1,3}$). \\ \hline
\textbf{Artifact} & \textbf{PANOPTES:} A reference architecture, multimodal RCA correlation engine, and IaC instantiation (Ansible/Helm) integrating OpenTelemetry, Prometheus, Loki, and Tempo. \\ \hline
\textbf{Evaluation Strategy} & Quantitative Controlled Experiment ($2\times2$ Factorial Design with paired, stratified fault campaigns) to validate Diagnostic Efficiency ($H_{1,1}$), Diagnostic Accuracy ($H_{1,2}$), and Bounded Resource Overhead ($H_{1,3}$). \\ \hline
\end{tabularx}
\par\medskip\ABNTEXfontereduzida\selectfont\textbf{Source:} Adapted from \citeonline{pimentel_2020} \par\medskip
\end{table}
